\section{Task 1: Multimodal Classification of Product Defects}
\label{sec:t1}
% Numbers: kode/results/task1/experiment_summary.json & per_class_f1_comparison.csv (python run_all.py --task 1).

We classify 16{,}000 test product units into seven defect classes (\texttt{normal}, \texttt{gores},
\texttt{noda}, \texttt{retak}, \texttt{gelembung}, \texttt{las\_dingin}, and \texttt{kontaminasi})
from 48{,}000 labelled training units. Each unit combines a $64 \times 64$ grayscale surface image
with a $3 \times 128$ press sensor recording (temperature, vibration, current). Evaluation on the
hidden test set uses macro-F1 over all seven classes. Our best model, a three-seed
\textbf{Cross-Attention} multimodal ensemble, achieves \textbf{0.8533} validation macro-F1
($0.8416 \pm 0.0023$ single-seed mean $\pm$ std).

\subsection{Data exploration}
\label{sec:t1-eda}

\paragraph{Class imbalance and metadata.} Nearly half of the 48{,}000 training units are
\texttt{normal} (23{,}888 units, 49.8\%), while the six defect classes range from 11.7\%
(\texttt{gores}, 5{,}623) down to 4.9\% (\texttt{kontaminasi}, 2{,}357); Table~\ref{tab:t1-dist}
lists the full distribution. Units come from 12 press machines, 150 raw-material batches, and three
balanced shifts (\texttt{pagi}, \texttt{siang}, \texttt{malam}).

\paragraph{Missing sensors and train--test domain shift.} In the training set, 2{,}352 units
(4.90\%) have completely missing sensor recordings (all \texttt{NaN}) and 3{,}863 units (8.05\%)
contain at least one \texttt{NaN}; the test set shows almost identical rates (813 units or 5.08\%
completely missing; 1{,}298 or 8.11\% with any \texttt{NaN}). In contrast, the production-line
composition shifts sharply between train and test: training data is dominated by Line~A
(35{,}964 units, 74.9\% vs.\ 12{,}036 or 25.1\% on Line~B), whereas the test set is dominated by
Line~B (9{,}607 units, 60.0\% vs.\ 6{,}393 or 40.0\% on Line~A).

\begin{table}[t]
  \centering
  \caption{Class distribution in the 48{,}000 training units and the 16{,}000 submitted test
  predictions (\texttt{prediksi\_klasifikasi.csv}).}
  \label{tab:t1-dist}
  \small
  \begin{tabular}{lrrrr}
    \toprule
    Class & Train count & Train share & Test pred.\ count & Test pred.\ share \\
    \midrule
    \texttt{normal}      & 23{,}888 & 49.8\% & 9{,}619 & 60.1\% \\
    \texttt{gores}       &  5{,}623 & 11.7\% & 1{,}240 &  7.8\% \\
    \texttt{noda}        &  4{,}773 &  9.9\% & 1{,}332 &  8.3\% \\
    \texttt{las\_dingin} &  4{,}323 &  9.0\% & 1{,}326 &  8.3\% \\
    \texttt{retak}       &  3{,}606 &  7.5\% &   806   &  5.0\% \\
    \texttt{gelembung}   &  3{,}430 &  7.1\% &   960   &  6.0\% \\
    \texttt{kontaminasi} &  2{,}357 &  4.9\% &   717   &  4.5\% \\
    \midrule
    Total                & 48{,}000 & 100.0\% & 16{,}000 & 100.0\% \\
    \bottomrule
  \end{tabular}
\end{table}

\subsection{Preprocessing and architectures}
\label{sec:t1-arch}

\paragraph{Preprocessing.} Training data is split 80/20 with class stratification
(\texttt{random\_state=42}; 38{,}400 train / 9{,}600 validation), and all normalization statistics
are fitted on the training split only. Because raw sensor arrays are stored in \texttt{float16},
valid (non-\texttt{NaN}) readings are cast to \texttt{float64} before computing per-channel means
and standard deviations; computing statistics directly in \texttt{float16} overflows on 4.9M entries
and zeroes out the sensor branch. Missing entries are filled with 0.0 after standardization. During
training we apply random horizontal/vertical flips to images and modality dropout ($p=0.15$, zeroing
the entire sensor tensor) so the model learns to fall back on the image when sensors are absent.

\paragraph{Encoders and fusion strategies.} Both modalities use compact residual
networks~\cite{he2016} producing 128-dimensional representations:
\begin{itemize}
  \item \textbf{Image encoder (2D ResNet):} initial $3\times 3$ convolution (16 channels) followed by
  four strided residual stages (32, 64, 128, 128 channels) reducing $64\times 64$ to $4\times 4$, then
  global average pooling and a linear projection to $\mathbb{R}^{128}$ (or 16 spatial tokens of dim
  128 for attention).
  \item \textbf{Signal encoder (1D ResNet):} initial kernel-5 convolution (16 channels) followed by
  four strided 1D residual stages (32, 64, 128, 128 channels) reducing length 128 to 8, then global
  average pooling and projection to $\mathbb{R}^{128}$ (or 8 temporal tokens for attention).
  \item \textbf{Late fusion:} concatenates the two 128-d embeddings and passes them through a
  two-layer MLP ($256 \to 128 \to 7$) with BatchNorm, ReLU, and dropout $0.3$.
  \item \textbf{Cross-attention fusion:} multi-head attention~\cite{vaswani2017} (4 heads,
  $d=128$) where the 16 image spatial tokens query the 8 signal temporal tokens, followed by residual
  LayerNorm, mean pooling of both streams, concatenation, and an MLP classifier.
\end{itemize}

\subsection{Modality value, label noise, domain shift, and missing sensors}
\label{sec:t1-experiments}

All comparison experiments train for 15 epochs with AdamW ($\text{lr}=10^{-3}$, weight decay
$10^{-4}$), cosine annealing, batch size 128, and seed 42; the final ensemble trains for 18 epochs
across three seeds (42, 123, 999). Before full training, a sanity check overfits one 32-sample batch
for 40 epochs (\texttt{train\_eval.py --sanity-check}), verifying cross-entropy loss $< 0.1$ and
macro-F1 $= 1.00$.

\begin{table}[t]
  \centering
  \caption{Per-class validation F1 and overall macro-F1 (9{,}600 validation units, seed 42).
  Visual defects rely on the image encoder, \texttt{las\_dingin} relies purely on the sensor
  encoder, and multimodal fusion excels on every class.}
  \label{tab:t1-modality}
  \small
  \begin{tabular}{lrrrr}
    \toprule
    Class & Image-Only & Sensor-Only & Late Fusion & Cross-Attention \\
    \midrule
    \texttt{normal}      & 0.827 & 0.704 & 0.899 & \textbf{0.905} \\
    \texttt{gores}       & 0.710 & 0.049 & 0.750 & \textbf{0.755} \\
    \texttt{noda}        & 0.872 & 0.092 & \textbf{0.874} & 0.871 \\
    \texttt{retak}       & 0.521 & 0.673 & 0.718 & \textbf{0.735} \\
    \texttt{gelembung}   & 0.556 & 0.760 & 0.776 & \textbf{0.794} \\
    \texttt{las\_dingin} & 0.000 & \textbf{0.936} & 0.933 & 0.934 \\
    \texttt{kontaminasi} & 0.884 & 0.008 & 0.899 & \textbf{0.899} \\
    \midrule
    \textbf{Macro-F1}    & 0.624 & 0.460 & 0.836 & \textbf{0.842} \\
    \bottomrule
  \end{tabular}
\end{table}

\paragraph{Value of each modality.} Table~\ref{tab:t1-modality} reveals a clean division between
modalities. The \textbf{Image-Only} 2D ResNet (macro-F1 0.624) detects surface defects accurately
(\texttt{kontaminasi} 0.884, \texttt{noda} 0.872, \texttt{gores} 0.710) but is completely blind to
\texttt{las\_dingin} (0.000). Conversely, the \textbf{Sensor-Only} 1D ResNet (macro-F1 0.460)
detects \texttt{las\_dingin} with 0.936 F1 and outperforms images on \texttt{gelembung} (0.760 vs.\
0.556) and \texttt{retak} (0.673 vs.\ 0.521), yet cannot see surface-only defects
(\texttt{kontaminasi} 0.008, \texttt{gores} 0.049, \texttt{noda} 0.092). Combining both modalities
adds over +0.21 macro-F1: \textbf{Late Fusion} reaches 0.8356 and \textbf{Cross-Attention} reaches
\textbf{0.8419}, with the largest cross-attention gains on the subtle dual-modality classes
\texttt{gelembung} (+0.018) and \texttt{retak} (+0.017).

\begin{table}[t]
  \centering
  \caption{Summary of Task 1 robustness, domain-shift, and final 3-seed validation results
  (9{,}600 validation units; from \texttt{kode/results/task1/experiment\_summary.json}).}
  \label{tab:t1-results}
  \small
  \begin{tabular}{lrrr}
    \toprule
    Experiment / Model & Overall F1 & Line A F1 & Line B F1 \\
    \midrule
    Image-Only 2D ResNet                                      & 0.6243 & --     & --     \\
    Sensor-Only 1D ResNet                                     & 0.4601 & --     & --     \\
    Late Fusion (global norm, LS $\epsilon=0.1$)              & 0.8356 & 0.8755 & 0.7004 \\
    Late Fusion (no label smoothing, $\epsilon=0.0$)          & 0.8369 & --     & --     \\
    Late Fusion + Per-Line Norm (image \& sensor)             & 0.8269 & 0.8704 & 0.6807 \\
    Cross-Attention (global norm, LS $\epsilon=0.1$)          & 0.8419 & --     & --     \\
    Cross-Attention + Per-Line Norm (image \& sensor)         & 0.8415 & 0.8797 & 0.7099 \\
    Cross-Attention with 100\% Missing Sensor                 & 0.5777 & --     & --     \\
    Cross-Attention on Natural Missing-Sensor Val Subset      & 0.5497 & --     & --     \\
    \midrule
    Final Cross-Attention (3 seeds: 42, 123, 999)             & $0.8416 \pm 0.0023$ & -- & -- \\
    \textbf{Final 3-Seed Probability Ensemble}                & \textbf{0.8533} & \textbf{0.8886} & \textbf{0.7328} \\
    \bottomrule
  \end{tabular}
\end{table}

\paragraph{Imperfect labels.} Using the trained multimodal model to flag high-confidence
disagreements (predicted probability $\ge 0.85$) identifies \textbf{339 likely mislabels} out of
9{,}600 validation units (3.53\%). Comparing standard cross-entropy ($\epsilon=0.0$, macro-F1
0.8369) against label smoothing~\cite{szegedy2016} ($\epsilon=0.1$, macro-F1 0.8356) yields nearly
identical scores on the noisy validation labels (difference 0.0013), while preventing overconfident
memorization of corrupted inspector labels; we keep $\epsilon=0.1$ for regularization in the final
training recipe.

\paragraph{Domain shift across production lines.} Evaluating by production line exposes a large
domain gap: under global normalization, Late Fusion scores 0.8755 on Line~A but drops to 0.7004 on
Line~B ($-0.175$). Standardizing both images and sensors per line slightly lowers Late Fusion
(0.8269 overall; Line~B 0.6807) because Line~B has $3\times$ fewer training samples for estimating
statistics, whereas \textbf{Cross-Attention + Per-Line Norm} lifts both lines (Line~A 0.8797, Line~B
0.7099; overall 0.8415). Because our automated selection criterion compared Late Fusion per-line
against global normalization before picking the architecture, the final ensemble used
Cross-Attention with global normalization; ensembling three seeds further lifts Line~B to
\textbf{0.7328} (Line~A \textbf{0.8886}).

\paragraph{Missing modality robustness.} Trained with 15\% sensor modality dropout, the chosen
Cross-Attention model retains \textbf{0.5777} macro-F1 even when 100\% of validation sensors are
wiped to \texttt{NaN} (approaching the dedicated Image-Only ceiling of 0.6243), and achieves
\textbf{0.5497} macro-F1 on the naturally missing-sensor validation subset---confirming graceful
fallback to visual features rather than collapsing.

\subsection{Final model and test predictions}
\label{sec:t1-final}

The final model trains \textbf{Multimodal Cross-Attention} for 18 epochs across three random seeds:
seed 42 (0.8428), seed 123 (0.8435), and seed 999 (0.8384), giving a mean validation macro-F1 of
$0.8416 \pm 0.0023$. Averaging the softmax probabilities of the three checkpoints raises validation
macro-F1 to \textbf{0.8533} (Line~A 0.8886, Line~B 0.7328). Applied to the 16{,}000 hidden test
units, \texttt{prediksi\_klasifikasi.csv} predicts all seven classes in proportions consistent with
the training distribution and the Line~A $\to$ Line~B shift (Table~\ref{tab:t1-dist}).
